% 1. Preamble
\documentclass[12pt, a4paper]{article}

% Import Packages
\usepackage[utf8]{inputenc}        % Encoding
\usepackage{array}
\usepackage{needspace}
\usepackage{longtable}
\usepackage{amsmath, amssymb}     % Advanced math formulas
\usepackage{mathtools}           % Math enhancements
\usepackage{graphicx}            % Image insertion
\usepackage{hyperref}            % Interactive hyperlinks & PDF bookmarks
\usepackage{color}               % Color text and backgrounds
\usepackage{geometry}            % Page geometry
\usepackage{comment}             % Multi-line comments
\usepackage{titling}             % Custom title formatting

\geometry{a4paper, margin=1in}

% Custom Command Definitions
\newcommand{\gmunu}{g_{\mu\nu}}
\newcommand{\Tmunu}{T_{\mu\nu}}
\providecommand{\norm}[1]{\left\lVert#1\right\rVert}

% Hyperref Setup for Clean PDF Links
\hypersetup{
    colorlinks=true,
    linkcolor=blue,
    urlcolor=blue,
    pdftitle={TITLE PLACEHOLDER: A Hypothesis on Space, Time and Gravity},
    pdfauthor={Sandeep Lakshminarayan Chiluveru}
}

% Metadata & Custom Title Banner
\pretitle{\begin{center}\hbox to\linewidth{\leaders\hbox{\normalfont\ttfamily\normalsize =}\hfill}\par\vspace{0.5em}\Huge}
\posttitle{\par\vspace{0.5em}\hbox to\linewidth{\leaders\hbox{\normalfont\ttfamily\normalsize =}\hfill}\end{center}}

% --- WORKING TITLE: replace once the hypotheses are fixed ---
\title{On the Elasticity and Origin of Spacetime Curvature: \\ Two Working Hypotheses on Gravity}
\author{Sandeep Lakshminarayan Chiluveru}
\date{\today}

\begin{document}

\maketitle

% =====================================================================
% ABSTRACT
% =====================================================================
\begin{abstract}
% DRAFT -- tighten once H1/H2 are formalized with equations.
General relativity identifies gravity with spacetime curvature:
mass-energy curves spacetime, and free bodies follow the resulting
geodesics. This paper sets out two working hypotheses that revisit
that identification from opposite directions. Hypothesis~1 proposes
that space itself is an elastic medium, and that the amount of
curvature a given mass produces is governed by that medium's
elasticity, so that the strength of gravity is tied directly to the
elastic response of space rather than treated as a bare coupling
constant. Hypothesis~2 reverses the usual causal order: it proposes
that curvature is not the source of gravity but a \emph{consequence}
of it, and that the gravitational force itself originates either from
a constant, universe-scale acceleration (cosmic expansion) or from a
constant rotation (spin) of space, with mass ``sinking'' into space
under that acceleration or spin and thereby bending it, in a manner
licensed, by hypothesis, through the equivalence principle. Both
hypotheses are stated here in their roughest form, as working ideas to
be formalized, checked against known limits, and tested against
observational constraints -- not as completed theories. This paper
does not yet claim to reproduce the quantitative successes of general
relativity (inverse-square attraction, perihelion precession, light
bending, gravitational-wave speed); establishing that recovery is the
first item of future work flagged throughout.
\end{abstract}

% =====================================================================
\section{Introduction}
% =====================================================================

\paragraph{Summary.}
% Plain-language statement of the question driving the paper and why
% the standard picture (GR + QFT) leaves it open. 2-4 sentences.
[PLACEHOLDER]

\paragraph{At a glance.}
% A short table naming each hypothesis, the question it answers, and
% where in the paper it is formalized -- mirrors the "six verbs" table
% convention. Fill in rows per hypothesis (H1, H2, H3, ...).
\begin{center}
\footnotesize
\begin{tabular}{@{}>{\raggedright\arraybackslash}p{1.6cm}>{\raggedright\arraybackslash}p{4.0cm}>{\raggedright\arraybackslash}p{5.2cm}>{\raggedright\arraybackslash}p{2.0cm}@{}}
\hline
Hypothesis & Question it answers & Claim, in one line & Where \\ \hline
H1 & What sets the amount of curvature for a given mass? & Gravity's strength is tied to the elasticity of space; bend amount $\propto$ elasticity & \S\ref{subsec:h1} \\
H2a & What causes curvature, and where does the force of gravity come from? & Curvature is a consequence, not the cause, of gravity; gravity is a constant acceleration from cosmic expansion, equivalent to a gravitational field by the equivalence principle & \S\ref{subsec:h2} \\
H2b & (alternative to H2a) & Gravity instead comes from a constant spin of space about an axis; the resulting centrifugal effect is equivalent to a gravitational field & \S\ref{subsec:h2} \\
H3 & Why do some black holes (supermassive) far exceed the mass stellar collapse can produce? & Space is not smooth; large tears/inconsistencies in its fabric, from uneven cosmic expansion and turbulence, form the largest black holes directly & \S\ref{subsec:h3} \\
H4 & What seeds the accretion that forms stars, planets, and asteroids? & Small, widespread tears/inconsistencies in space trap matter and seed accretion; every body in space has a black hole at its center & \S\ref{subsec:h4} \\
H5 & What, specifically, are supermassive black holes, if not grown stellar remnants? & Supermassive black holes are large vortices formed where oppositely-expanding regions of space collide; they are not, in general, of stellar-collapse origin & \S\ref{subsec:h5} \\
H6 & Why do stars at a galaxy's edge orbit faster than visible mass alone predicts? & A galaxy is matter settled into the depression around a central vortex (the SMBH); the vortex's strength and reach set the rotation curve directly, so dark matter is not needed & \S\ref{subsec:h6} \\ \hline
\end{tabular}
\end{center}

\paragraph{Background and motivation.}
% Longer version of Summary: what GR gets right (equivalence principle,
% geodesic motion, confirmed predictions -- lensing, perihelion precession,
% gravitational waves, time dilation/GPS), what it leaves unexplained
% (singularities, quantization of gravity, dark energy/matter as currently
% framed, the measurement/arrow-of-time problem if relevant), and where
% this paper's hypotheses are meant to sit relative to that frontier.
% Be explicit about which open problem each hypothesis targets.
[PLACEHOLDER]

\begin{quote}
\textbf{Thesis.}
\emph{[One or two sentences stating the paper's organizing idea at the
level of a thesis statement, the way a thesis is stated before the
formal machinery arrives. State what kind of claim this is --
a reinterpretation, a new mechanism, an extension, an analogy under
test -- so the reader knows how much weight it is meant to carry.]}
\end{quote}

\paragraph{Scope and disclaimers.}
% State plainly what this is NOT: not a replacement for GR unless and
% until it reproduces GR's confirmed results in the appropriate limit;
% not a claim of experimental confirmation; a hypothesis awaiting
% derivation and test. Name the limit(s) in which it must reduce to
% known physics (e.g. weak-field limit, low-velocity limit, Newtonian
% limit) -- this is the single most important falsifiability anchor.
[PLACEHOLDER]

% =====================================================================
\section{Foundational Basis and Prior Work}
\label{sec:foundations}
% =====================================================================
% For EACH hypothesis, this section should eventually contain a short,
% honest assessment of how strongly it rests on established physics.
% Suggested sub-structure per hypothesis (duplicate as needed):

\subsection{Established results this paper relies on}
\begin{itemize}
  \item \textbf{Equivalence principle} (Einstein, 1907/1911; tested to
    high precision by E\"otv\"os-type and modern torsion-balance
    experiments): locally and instantaneously, a uniformly accelerated
    frame is indistinguishable from a uniform gravitational field. H2
    uses this unchanged, but extends its scope from a local,
    instantaneous statement to a claim about the global \emph{origin}
    of gravity -- an extension, not a restatement.
  \item \textbf{Einstein field equations}, $G_{\mu\nu} = 8\pi G\,T_{\mu\nu}$:
    curvature is set by mass-energy through the fixed constant $G$. H1
    keeps this structure but reinterprets $G$ (or the coupling it
    plays) as a measurable elasticity rather than a bare constant --
    a reinterpretation of an existing equation, not a new one, until
    ``elasticity'' is given independent empirical content.
  \item \textbf{Measured constancy of $G$}: Newton's constant is
    measured to be constant across laboratory tests to good precision.
    Any version of H1 in which elasticity varies must explain why that
    variation has not already shown up in existing $G$ measurements.
  \item \textbf{Cosmic expansion and its acceleration} (Hubble's law;
    the accelerating expansion inferred from Type Ia supernovae,
    Riess et al. 1998, Perlmutter et al. 1999): real, but measured to
    be extremely small and essentially uniform on the relevant scales
    -- not currently known to be locally concentrated around
    individual masses the way H2a would need.
  \item \textbf{Observed isotropy of the cosmic microwave background}:
    used to place very tight upper bounds on any global rotation of
    the universe, directly bearing on H2b.
  \item \textbf{Smoothness (differentiability) of the spacetime
    manifold}: a foundational assumption of classical GR; a literal
    topology change (a ``tear'') is excluded by it. Bears directly on
    H3/H4.
  \item \textbf{Mass bound on stellar-collapse black holes}
    (progenitor mass limits; the pair-instability mass gap, roughly
    $50$--$150\,M_\odot$): well established, and correctly used as the
    starting observation for H3.
  \item \textbf{Observed supermassive black holes at high redshift}:
    the ``seed problem'' -- a genuine, open question H3 correctly
    identifies, already addressed in the literature by direct-collapse,
    heavy-seed, and super-Eddington-accretion proposals.
  \item \textbf{Helioseismology, terrestrial seismology, and
    planetary geodesy}: direct observational probes of the interiors
    of the Sun, Earth, and other bodies, bearing directly on H4's
    claim that a black hole sits at the center of every body.
  \item \textbf{Quantum foam} (Wheeler, 1955) and \textbf{primordial
    black holes} (Zel'dovich \& Novikov, 1966; Hawking, 1971): existing,
    legitimate research programs nearest in spirit to H3/H4's
    ``inconsistencies in space,'' discussed below.
  \item \textbf{Homogeneity and isotropy of cosmic expansion (the FLRW
    model)}: supported by CMB isotropy and large-scale-structure
    surveys; the expansion rate is a single scale factor, not a
    directional flow field with distinct colliding regions. Bears
    directly on H5.
  \item \textbf{Black hole rotation, frame dragging, and
    analogue-gravity vortex experiments} (Kerr metric; Lense--Thirring
    effect, now measured via accretion-disk spin diagnostics; Unruh,
    1981, and later draining-fluid-vortex demonstrations of
    superradiance): the genuine, if indirect, point of contact for
    H5's vortex picture.
  \item \textbf{Galaxy rotation curves} (Rubin \& Ford, 1970s onward):
    the original, and still central, dark-matter evidence H6 targets.
  \item \textbf{Independent dark-matter evidence beyond rotation
    curves}: the Bullet Cluster's lensing/gas offset; the relative
    heights of the CMB acoustic peaks; the growth rate of large-scale
    structure. Each is logically independent of rotation curves and
    must be addressed by any replacement, not just the curves
    themselves.
  \item \textbf{Modified Newtonian Dynamics (MOND)} (Milgrom, 1983):
    the existing, well-studied alternative to dark matter closest in
    spirit to H6, with known successes (individual rotation curves)
    and known difficulties (cluster scales, the Bullet Cluster, the
    CMB).
\end{itemize}

\subsection{Where each hypothesis is standard, extrapolated, or speculative}
\textbf{H1 (elasticity of space).} The claim that mass-energy curves
space, and that the amount of curvature follows from a fixed coupling,
is (a) standard GR. The claim that this coupling should be understood
as a physical \emph{elasticity} is (b) an extrapolation: it has real
precedent as an analogy (spacetime-as-elastic-medium pictures appear
in induced-gravity and analogue-gravity literature, \S\ref{sec:foundations}.3
below) but is not itself an established, independently measured
quantity. As currently stated -- with no independent definition of
elasticity and no claimed deviation from $G_{\mu\nu}=8\pi G T_{\mu\nu}$
-- H1 is (c) not yet distinguishable from a relabeling of $G$, which
makes it unfalsifiable in its present form. It becomes a genuine,
moderate-strength hypothesis exactly to the extent that elasticity is
defined as something that could, in principle, be measured
\emph{independently} of $G$ and found to vary.

\textbf{H2 (gravity sourced by expansion or spin).} The equivalence
principle itself is (a) standard and extremely well tested. Using it
to say that \emph{a} sustained acceleration feels like gravity is
correct. Using it to say that cosmic expansion or cosmic rotation
\emph{is the acceleration that constitutes all gravity, everywhere,
including the mass-centered field around the Earth} is (c) a new
postulate with no direct precedent and, as detailed in
\S\ref{subsec:h2}, currently in tension with the measured size and
uniformity of cosmic expansion, with CMB-isotropy bounds on cosmic
rotation, and with the basic mass-centered, inverse-square shape of
every gravitational field actually observed. This is the weakest part
of either hypothesis in its current form, precisely because it is the
most novel claim and the one most directly checked against existing,
precise observations.

\textbf{H3 (tears as the origin of supermassive black holes).} That
stellar-collapse black holes are mass-limited, and that supermassive
black holes pose an unsolved seed problem, is (a) standard and
accurate. That the resolution is a literal tear in spacetime's fabric
is (c) a new postulate: GR's manifold is smooth by construction, so
this is not an extrapolation of an existing pattern but a proposed
exception to a foundational assumption, and it currently has to
compete with, rather than build on, the mainstream candidate answers
(direct collapse, heavy seeds, super-Eddington growth) without yet
stating what distinguishes it from them observationally.

\textbf{H4 (small inconsistencies seeding all accretion; a black hole
at the center of every body).} That accretion builds stars, planets,
and asteroids, and that growing mass deepens a gravitational well, is
(a) standard accretion physics, just stated in this paper's own
vocabulary. That the \emph{trigger} for accretion is a small tear or
hole in space, rather than ordinary gravitational instability in a
gas or dust cloud, is (c) a new, currently unnecessary postulate --
the mainstream mechanism already accounts for the observation without
it. The claim that every body has a black hole at its center is also
(c), and is the most strongly constrained part of either hypothesis:
it is checked directly against existing interior-structure data
(helioseismology, seismology, geodesy) and, as stated without
qualification, does not currently match it.

\textbf{H5 (supermassive black holes as colliding-space vortices).}
The starting observation -- that a single star's collapse cannot
directly produce a supermassive black hole -- is (a) standard and
correctly stated. The vortex/rotation imagery has (b) a real,
extrapolatable point of contact: real black holes do drag spacetime
when they rotate, and that rotational mathematics is exactly what
analogue-gravity vortex experiments exploit. But the specific sourcing
mechanism -- large regions of space expanding in opposite directions
and colliding -- is (c) a new postulate with no observational
counterpart in the standard, homogeneous-and-isotropic picture of
cosmic expansion, and it has to explain away, not just outpace, the
mainstream seed-growth picture (stellar-mass seeds plus mergers and
sustained accretion), which already accounts for supermassive black
holes' observed mass and spin without requiring colliding expansion.

\textbf{H6 (galaxies as vortex depressions, replacing dark matter).}
The target observation -- flat rotation curves -- is (a) standard and
correctly identified. Attributing it to the central vortex (the
galaxy's supermassive black hole) is (c) a new postulate that, as it
stands, faces a direct scale problem: the black hole's gravitational
reach is three to four orders of magnitude smaller than the radii
where flat rotation curves are actually measured, so the mechanism as
stated does not yet reach the regime it is meant to explain. It is
also (c) silent on the independent, non-rotation-curve evidence for
dark matter (Bullet Cluster, CMB acoustic peaks, structure growth),
each of which would need its own account. This is the most
empirically consequential hypothesis in the set so far, precisely
because dark matter is supported by several independent lines of
evidence rather than one -- a replacement has to clear all of them,
not just the one it was built to match.

\subsection{Relation to existing programs}
\textbf{H1} sits closest to \emph{induced gravity} (Sakharov, 1967),
where $G$ itself is proposed to emerge from vacuum fluctuations rather
than being fundamental, and to various \emph{analogue-gravity}
treatments that model spacetime's response to mass with
elastic-medium mathematics as a teaching or effective-theory device.
Verlinde's entropic-gravity proposal (2011) is a different but related
example of treating gravity as emergent from a coarser underlying
layer rather than fundamental. None of these establish that spacetime
literally possesses a material elastic modulus; they are precedents
for the \emph{style} of H1's claim, not confirmations of it.

\textbf{H2's} ``sinking into space'' picture has a rigorous existing
counterpart worth knowing about: the \emph{river model of gravity}
(Painlev\'e, 1921; Gullstrand, 1922; popularized by Hamilton \& Lisle,
2008), in which space itself is pictured as flowing inward toward a
mass -- mathematically just a coordinate choice (Painlev\'e--Gullstrand
coordinates) fully equivalent to standard Schwarzschild/GR, not a
modification of it. That model gives ``mass makes space flow/sink''
a precise, already-solved form, but notably sources the flow from the
\emph{local} mass, not from cosmic expansion or universal spin -- the
opposite of what H2a/H2b propose. A genuine, non-trivial rotating
universe is also an existing GR solution (G\"odel, 1949), but it is
non-expanding, permits closed timelike curves, and is not considered a
model of our universe; current observations bound our universe's
actual rotation far below what H2b would require. Distinguishing
H2 from the river model -- i.e. stating clearly whether the ``sinking''
is sourced locally (as in the river model) or globally (as currently
written) -- is the single most useful next step for sharpening H2.

\textbf{H3 and H4} sit closest to two existing, separate programs that
this paper's single notion of ``inconsistency'' currently blends
together. \emph{Quantum foam} (Wheeler, 1955) proposes that spacetime
is topologically non-trivial at the Planck scale due to quantum
fluctuations -- a real precedent for ``space is not smooth,'' but at a
scale enormously below anything relevant to black-hole formation, with
no established bridge to that scale. \emph{Primordial black holes}
(Zel'dovich \& Novikov, 1966; Hawking, 1971) propose that density
fluctuations in the early universe collapse directly into black holes
of a wide range of masses, independent of stellar collapse, and have
been studied both as dark-matter candidates and as possible seeds for
structure formation -- a real precedent for ``small holes seed
accretion'' and even for ``some black holes are not stellar remnants,''
but sourced by early-universe density fluctuations, not by tears from
cosmic expansion or turbulence. Reframing H3/H4 explicitly in terms of
one of these two existing programs -- rather than the new term
``inconsistency,'' which currently has no agreed definition in either
GR or quantum gravity -- would let the hypotheses inherit an existing,
worked-out mathematical machinery instead of having to build one from
scratch.

\textbf{H5} sits closest to two existing, separate programs as well.
\emph{Analogue gravity} (Unruh, 1981, and subsequent draining-vortex
experiments) shows that a real fluid vortex can mathematically mimic
a rotating black hole's horizon and ergosphere for waves travelling
in that fluid -- a precedent for the mathematics of a ``vortex,'' but
one built because sound waves in a draining fluid obey the same wave
equation as fields near a Kerr horizon, not because spacetime itself
is a fluid that forms vortices. \emph{Bubble-collision black-hole
formation} during an early-universe phase transition (studied in
certain inflationary and phase-transition cosmological models) is the
closest precedent for ``a collision forms a black hole,'' but it is a
mechanism specific to the very early universe and to a particular
particle-physics scenario, not a general or ongoing process tied to
cosmic expansion as such. As with H3/H4, stating which of these two
programs H5 means to extend -- and on what grounds the collision
mechanism applies well beyond the early-universe regime where its
nearest precedent is actually set -- is the clearest way to sharpen it.

\textbf{H6} sits closest to \emph{MOND} and its relativistic
extensions (TeVeS and successors), which already pursue exactly H6's
stated goal -- explaining rotation curves without a new matter
species -- by modifying the force law or the field equations at low
accelerations, rather than by invoking a vortex. MOND's real track
record is informative either way: it fits individual galaxy rotation
curves remarkably well, which is evidence that \emph{some}
modification-of-gravity approach can work at that scale: but it
has never comfortably accounted for the Bullet Cluster or the CMB
acoustic-peak ratios without reintroducing some form of unseen mass.
If H6 is pursued further, inheriting MOND's success on rotation
curves is the easy part; it will need its own answer to the harder
part MOND has not solved, rather than treating rotation curves as
the whole problem.

% =====================================================================
\section{Notation and Primitive Objects}
\label{sec:notation}
% =====================================================================
% Define, once and for all, the objects the rest of the paper uses:
% the manifold, metric signature convention, coordinates, what counts
% as an "event," "worldline," "observer," and any new object this
% paper introduces (mirror the atom/core style of definition: state
% the object, its components, and its domain, precisely).

\begin{itemize}
  \item Manifold / spacetime: [PLACEHOLDER, e.g. a 4-dimensional
    Lorentzian manifold $(M, \gmunu)$ with signature $(-,+,+,+)$]
  \item Units and conventions: [PLACEHOLDER, e.g. $c=1$ unless stated,
    $G$ retained explicitly]
\end{itemize}

\subsection{Core Terms: Elasticity and Fluidity of Space}
\label{subsec:core_terms}
Hypotheses H1--H6 all lean on two intuitive properties of space --
its \emph{elasticity} (how much it resists bending) and its
\emph{fluidity} (how it flows, expands, spins, or shears locally).
Rather than inventing new mathematics for these, each is defined here
as a generalization of an object GR already has, so that ``standard
GR'' is the special case where the generalization is held fixed, and
each hypothesis's new content is exactly the claim that it isn't.

\paragraph{Elasticity: the compliance field $\Sigma$.}
Standard GR writes $G_{\mu\nu} = (8\pi G/c^4)\,T_{\mu\nu}$, a fixed
coupling. The combination $c^4/(8\pi G) \approx 4.98\times10^{42}\,
\mathrm{N}$ is already informally called the \emph{stiffness of
spacetime} (Thorne and others) -- it is how much stress-energy it
takes to produce a given amount of curvature, i.e. exactly the
quantity H1 calls elasticity, already present in the theory as a
constant. Define
\begin{equation}
\Sigma(x) \;:=\; \frac{c^4}{8\pi G(x)}, \qquad
G_{\mu\nu} \;=\; \frac{1}{\Sigma(x)}\,T_{\mu\nu}.
\end{equation}
Standard GR is the case $\Sigma(x) = \Sigma_0$, constant everywhere.
\textbf{H1's content is exactly, and only, the claim that $\Sigma$
varies} -- with density, curvature, scale, or something else to be
specified -- subject to reducing to the measured constant $\Sigma_0$
everywhere $G$ has actually been tested (\S\ref{subsec:h1}). Writing
$\Sigma(x)$ rather than ``elasticity'' turns H1 into a statement about
one specific field, with one specific job: say what $\Sigma$ depends
on, and by how much, and the hypothesis becomes checkable.

\paragraph{Fluidity: the kinematic triplet $(\theta, \sigma_{\mu\nu},
\omega_{\mu\nu})$.}
GR already has a standard, rigorous way to describe a ``flow'' of
space: take a congruence of worldlines (e.g. the worldlines of
observers at rest with respect to local matter) and decompose how
nearby worldlines move relative to one another into three pieces --
\begin{itemize}
  \item $\theta$, the \textbf{expansion scalar}: the local
    rate at which a small volume of worldlines grows or shrinks.
    This is precisely cosmic expansion, generalized to be local: in
    the standard cosmological model $\theta = 3H$, three times the
    Hubble parameter. \textbf{H2a's ``expanding space''} is a
    statement about $\theta$.
  \item $\omega_{\mu\nu}$, the \textbf{vorticity (rotation) tensor}:
    the local rotation of nearby worldlines about each other.
    \textbf{H2b's ``spin of space''} and the rotational character of
    \textbf{H5's and H6's ``vortex''} are statements about
    $\omega_{\mu\nu}$ (or its vector form $\omega^\mu$).
  \item $\sigma_{\mu\nu}$, the \textbf{shear tensor}: local
    stretching or squeezing without a change in volume or net
    rotation -- distortion into an ellipsoid rather than a sphere.
    This is the natural home for \textbf{H3/H4/H6's ``turbulence,''
    ``tears,'' ``cracks,'' and ``channels''}: a shear-dominated region
    stretches matter preferentially along one direction, which is
    exactly the qualitative picture of a filament or a crack, as
    opposed to $\theta$ (uniform swelling) or $\omega_{\mu\nu}$
    (rotation).
\end{itemize}
These three are not independent of curvature or of matter: they are
tied together by the \textbf{Raychaudhuri equation},
\begin{equation}
\frac{d\theta}{d\lambda} = -\frac{1}{3}\theta^2 -
\sigma_{\mu\nu}\sigma^{\mu\nu} + \omega_{\mu\nu}\omega^{\mu\nu} -
R_{\mu\nu}u^\mu u^\nu,
\end{equation}
where the last term is fixed, through the Einstein field equations,
by the local mass-energy density. This is the precise, already-
established sense in which ``fluidity'' cannot float free of a
mass-energy source (the open requirement raised against H2, H4, H5,
and H6 throughout \S\ref{sec:foundations} and \S\ref{sec:hypotheses}):
$\theta$, $\sigma_{\mu\nu}$, and $\omega_{\mu\nu}$ are properties of
how matter moves and curves spacetime, not properties spacetime can
be assigned independently of matter. \textbf{Using this formalism
does not rescue any hypothesis from that requirement} -- it hands you
the exact equation that requirement is written in, so that ``what
sources this'' becomes a specific term to compute rather than an open
question to wave at.

\paragraph{Inconsistency: the irregularity scalar $\mathcal{I}$.}
H3--H6 all invoke a third core term -- that space is not smooth, that
it is ``riddled with inconsistencies.'' This, too, already has a
rigorous counterpart to generalize rather than a word to leave
informal: standard cosmological \emph{perturbation theory} already
describes the real universe as a smooth background (homogeneous,
isotropic FLRW) plus deviations from it, and measures those
deviations with a dimensionless quantity (conventionally $\delta$,
the density contrast) relative to that background. Define, in the
same spirit, an \textbf{irregularity scalar}
\begin{equation}
\mathcal{I}(x) \;:=\; \ell^2\sqrt{\,(\theta(x)-\bar\theta)^2 \;+\;
2\,\sigma_{\mu\nu}(x)\sigma^{\mu\nu}(x) \;+\; 2\,\omega_{\mu\nu}(x)
\omega^{\mu\nu}(x)\,}\,,
\end{equation}
where $\bar\theta$ is the smooth background expansion rate at that
cosmic time and $\ell$ is a reference length scale fixed by
convention (so $\mathcal{I}$ is dimensionless). $\mathcal{I}(x)=0$
recovers perfectly smooth FLRW space; $\mathcal{I}(x)>0$ is a
``spacetime inconsistency'' in H3--H6's sense, with its
\emph{magnitude} setting how large a departure it is and its
\emph{dominant term} ($\theta$, $\sigma_{\mu\nu}$, or
$\omega_{\mu\nu}$) setting which taxonomy species it looks like
(\S\ref{subsec:taxonomy}). This is deliberately \emph{not} how the
Sphere species (an actual black hole) should be defined, however: GR
already has an exact, independent, and far better-established
criterion for that -- a \textbf{trapped surface} (Penrose, 1965),
a closed surface from which even outward-directed light cannot
escape -- and the formation of a horizon should be checked against
that criterion, not against a large value of $\mathcal{I}$. Treating
$\mathcal{I}$ as graded, continuous ``inconsistency'' for the
Vortex/Depression/Crack species, while keeping the Sphere species
tied to the trapped-surface criterion it already has, is what keeps
the taxonomy from quietly promoting an ordinary density fluctuation
into a black hole merely because some scalar got large.

\textbf{What remains open:} the threshold value of $\mathcal{I}$ at
which a region counts as a named species at all (as opposed to an
ordinary, already-expected density fluctuation, which also has
$\mathcal{I}>0$ in this definition), and the reference scale $\ell$,
are not fixed by anything stated so far -- exactly like $\Sigma$'s
dependence and the Raychaudhuri source term, this is a free choice
the hypotheses need to make explicitly before $\mathcal{I}$ does
predictive work rather than descriptive work.

\paragraph{How this maps onto the taxonomy (\S\ref{subsec:taxonomy}).}
A natural first pass: \emph{Sphere} = a trapped surface is present,
regardless of $\mathcal{I}$; \emph{Vortex} = $\mathcal{I}>0$,
$\omega_{\mu\nu}$-dominated, no trapped surface; \emph{Depression} =
$\mathcal{I}>0$, small and localized, $\theta<0$-dominated, no
trapped surface; \emph{Crack/channel} = $\mathcal{I}>0$,
$\sigma_{\mu\nu}$-dominated, elongated, no trapped surface. This is
offered as a starting point for filling in the taxonomy table's open
columns, not as a settled result.

% =====================================================================
\section{Hypotheses}
\label{sec:hypotheses}
% =====================================================================
% State each hypothesis as a single, precise, falsifiable sentence,
% followed by its formal statement and its motivation. Keep hypotheses
% independent of each other where possible, so a reader can accept one
% and reject another.

\subsection{Hypothesis 1: Gravity as Elasticity of Space}
\label{subsec:h1}
\textbf{Statement (rough form).} Space can bend, and space has
elasticity. The amount of bend produced by a given mass is set by
that elasticity. Gravity is therefore directly proportional to the
elasticity of space.

\textbf{Formal statement.}
% PLACEHOLDER -- needs a precise definition of "elasticity" (a scalar?
% a tensor field? a constant, or a field that varies with position,
% density, or curvature itself?) before an equation can be written.
% Sketch, to be replaced:
\begin{equation}
\text{(bend / curvature)} \;=\; f(\text{elasticity of space}) \times \text{(mass)}, \qquad f \text{ unspecified}
\end{equation}
Compare to the Einstein field equations, which already relate
curvature to mass-energy through a fixed coupling:
\begin{equation}
G_{\mu\nu} \;=\; 8\pi G\, T_{\mu\nu}.
\end{equation}

\textbf{Motivation.} Gravity in general relativity already states
that curvature is produced by mass-energy, through the constant $G$.
H1's new content, if any, is the claim that $G$ -- or whatever plays
its role -- is not a bare constant but the measurable elastic
response of a medium, in the way a material's stiffness sets how much
it deforms under a given load.

\textbf{What would falsify it / what is still undefined.}
(i) ``Elasticity'' is not yet defined as a measurable quantity --
until it is, the hypothesis cannot be distinguished from simply
renaming $G$. (ii) If elasticity is a global constant, H1 makes no
prediction that differs from standard GR and is not yet falsifiable;
it becomes a hypothesis with content only if elasticity is allowed to
vary (with density, curvature, or scale) in a way GR's constant $G$
does not. (iii) Any such variation must still reduce to Newtonian
gravity with the measured value of $G$ in the weak-field, laboratory
regime, where $G$ is measured to be constant to high precision.

\subsection{Hypothesis 2: Curvature as a Consequence, not the Cause, of Gravity}
\label{subsec:h2}
\textbf{Statement (rough form).} Bent space is a consequence of
gravity, not its cause. The force of gravity itself originates from
either (2a) a constant, ongoing acceleration of an expanding universe,
or (2b) a constant spin of space about an axis. By the equivalence
principle, a sustained acceleration of $9.81\,\mathrm{m/s^2}$ (an
elevator, a rocket) is indistinguishable from a gravitational field.
Mass ``sinks into'' space under this acceleration or spin, and that
sinking is what bends space; the amount of bend depends on the mass,
the elasticity of space (H1), and the speed of the expansion (2a) or
spin (2b).

\textbf{Formal statement.}
% PLACEHOLDER -- "sinking into space" needs a precise kinematic
% definition before this can be written as an equation. Sketch:
\begin{equation}
\text{(2a)} \quad \vec a_{\text{gravity}} \;\overset{?}{=}\; \vec a_{\text{cosmic expansion}}(\text{mass, position}), \qquad
\text{(2b)} \quad \vec a_{\text{gravity}} \;\overset{?}{=}\; \omega^2 r \,\hat\rho \ \text{(centrifugal form)}
\end{equation}

\textbf{Motivation.} The equivalence principle is real and
well-tested: locally, uniform acceleration and a uniform gravitational
field are experimentally indistinguishable. H2 extrapolates this local
equivalence into a global claim about the \emph{origin} of every
gravitational field.

\textbf{What would falsify it / open problems.}
(i) Known gravitational fields (Earth's, the Sun's) are centered on
specific masses, static to high precision in the lab frame, and fall
off as $1/r^2$ from that mass; the hypothesis as stated does not yet
show how a universe-wide expansion rate or a universe-wide spin rate
produces a field with exactly that mass-centered, inverse-square
structure rather than a uniform or axis-centered one. (ii) The
measured acceleration of cosmic expansion (from dark energy) is many
orders of magnitude too small, and far too uniform/non-local, to
supply $9.81\,\mathrm{m/s^2}$ at Earth's surface without an additional,
yet-unstated amplification mechanism tied to local mass. (iii) A
global rotation of the universe is tightly constrained by the observed
isotropy of the cosmic microwave background to be far too small to
source planetary-scale gravity; (2b) needs to say whether ``spin of
space'' means this cosmic rotation or a different, as yet undefined,
local spin. (iv) Both variants currently need the equivalence
principle to do more work than it is usually read as doing: it
licenses treating acceleration \emph{as} a gravitational field
locally, but does not by itself explain why mass generates an
inverse-square field centered on itself.

\subsection{Hypothesis 3: Large-Scale Inconsistencies (Tears) in Space as the Origin of the Largest Black Holes}
\label{subsec:h3}
\textbf{Statement (rough form).} Space is not smooth; it is riddled
with inconsistencies in its fabric. Black holes are one type of such
inconsistency. Stellar-core-collapse black holes cannot grow very
massive. The largest black holes instead arise directly from large
tears or gaps in space, produced by uneven cosmic expansion and
turbulence.

\textbf{Formal statement.}
% PLACEHOLDER -- "tear," "gap," and "inconsistency" are not yet
% geometric objects. A tear/topology-change is a specific, unusual
% claim in GR and needs to be defined before an equation is possible.
[PLACEHOLDER: define ``inconsistency'' as a geometric object --
a curvature singularity, a topological defect, a region where the
manifold is non-differentiable -- before this can be written formally]

\textbf{Motivation.} Correctly notes a real, open astrophysical
puzzle: stellar-collapse black holes are bounded in mass (by
progenitor mass and pair-instability effects, roughly tens to a few
hundred solar masses), yet supermassive black holes
($10^6$--$10^{10}\,M_\odot$) are observed already in place within the
first billion years after the Big Bang -- the ``seed problem'' for
supermassive black holes, a genuinely unsolved question in current
astrophysics.

\textbf{What would falsify it / open problems.}
(i) In standard GR, spacetime is a smooth (differentiable) manifold
by construction; a literal ``tear'' is a topology change, which is
classically forbidden and, even in candidate quantum-gravity
treatments, is a highly specialized, unresolved research question --
not an established mechanism this hypothesis can yet borrow. (ii) The
nearest legitimate existing idea -- that spacetime is not perfectly
smooth at the Planck scale (``quantum foam,'' Wheeler 1955) -- operates
at $\sim10^{-35}\,\mathrm{m}$, some 30+ orders of magnitude below
anything that could seed a supermassive black hole; no established
theory currently bridges that scale gap, so it is not yet available as
support. (iii) Mainstream astrophysics already has several candidate
answers to the seed problem this hypothesis is competing with --
direct-collapse black holes from massive gas clouds, heavy seeds from
dense early star clusters, and super-Eddington accretion -- and H3
needs to say what it explains that these do not, or what observation
would distinguish ``a tear in space'' from a direct-collapse seed.

\subsection{Hypothesis 4: Widespread Small Inconsistencies as the Seed of All Accretion}
\label{subsec:h4}
\textbf{Statement (rough form).} Inconsistencies in space are common
and widespread, not just the rare large ones behind H3. These smaller
inconsistencies are small tears or holes that trap particles, seeding
the accretion of matter into stars, planets, and asteroids. Accretion
itself then bends space further, accelerating further collection.
It would not be wrong to say a black hole lies at the center of every
body in space.

\textbf{Formal statement.}
[PLACEHOLDER: as with H3, ``trap'' and ``hole'' need a precise
geometric or field-theoretic definition before this is an equation
rather than a picture]

\textbf{Motivation.} Correctly identifies that accretion is how stars,
planets, and asteroids form, and that growing mass deepens a
gravitational well, which is consistent with standard accretion
theory (gravitational instability and core-accretion models for star
and planet formation). The closest existing idea to ``a small hole
seeding accretion'' is the primordial black hole (PBH) hypothesis
(Zel'dovich \& Novikov, 1966; Hawking, 1971), in which density
fluctuations in the very early universe collapse directly into small
black holes, which have since been proposed as possible seeds for
structure formation and as a dark-matter candidate.

\textbf{What would falsify it / open problems.}
(i) Standard star and planet formation is already well modeled by
gravitational self-collapse of gas and dust under ordinary Newtonian/
GR gravity, with no need for an additional ``trapping'' mechanism; H4
needs to say what observation it predicts that the standard
core-accretion picture does not. (ii) The claim that \emph{every} body
in space -- every planet, asteroid, and ordinary star -- has a black
hole at its center is directly checked against existing evidence and,
as stated, is in tension with it: the Sun's interior structure is
measured via helioseismology, Earth's via seismology and its geodetic
and magnetic-field behavior (which requires an ordinary convecting
liquid-metal core), and none show the signatures a central black hole
would produce (anomalous mass-radius relation, an accretion luminosity
or steady mass growth, or the tidal/precision-tracking signatures a
compact central mass would leave). This is the most directly
falsifiable claim among all four hypotheses, and currently the one
with the most existing evidence against it as stated; narrowing it --
e.g., to large bodies only, or to the very early universe before
ordinary structure formed, closer to how the actual PBH hypothesis is
framed -- would bring it back into contact with open questions rather
than settled observations.

\subsection{Hypothesis 5: Supermassive Black Holes as Vortices from Colliding, Oppositely-Expanding Regions of Space}
\label{subsec:h5}
\textbf{Statement (rough form).} Ordinary (stellar-scale) black holes
are minor turbulence, tears, or inconsistencies in space. Supermassive
black holes are categorically different: they form where large
regions of space, expanding in opposing directions, collide and
create a large vortex. Supermassive black holes are not, in general,
of stellar-collapse origin -- though that origin is not impossible in
principle -- and share more in common with cyclones and hurricanes
than with collapsed stars.

\textbf{Formal statement.}
[PLACEHOLDER: ``regions of space expanding in opposing directions''
and ``collision'' need to be defined relative to the standard
cosmological picture, in which the expansion rate (the Hubble
parameter) is a single, roughly uniform, non-directional quantity at
a given cosmic time -- not a vector flow field with distinct, opposing
currents]

\textbf{Motivation.} Correctly inherits the real seed problem from H3:
a single star's collapse cannot directly produce a $10^6$--$10^{10}
\,M_\odot$ object, so some categorically different formation channel
for the largest black holes is a reasonable thing to look for. The
vortex picture also has a genuine, if indirect, point of contact with
real physics: real (rotating, Kerr) black holes drag spacetime around
them (frame dragging / the Lense--Thirring effect, now measured via
accretion-disk precession and X-ray spectroscopy), and this rotational
behavior is precisely what motivates \emph{analogue-gravity}
experiments, in which a real draining-fluid vortex is built in the lab
because the mathematics of waves near its drain matches the
mathematics of waves near a rotating black hole's horizon (Unruh,
1981; demonstrated with real draining-bathtub vortices showing
rotational superradiance). That match is at the level of the wave
equations, however, not a claim that spacetime itself is a literal
fluid.

\textbf{What would falsify it / open problems.}
(i) The standard cosmological model (the homogeneous, isotropic FLRW
spacetime, supported by the measured isotropy of the cosmic microwave
background and of large-scale structure) describes expansion as a
single scale factor increasing over time, not as a flow field with
distinct regions expanding in different directions that could
collide; H5 needs to state what observation would reveal such
oppositely-expanding regions, since none has been reported. (ii) Even
granting a collision, converting the kinetic energy of ``colliding
space'' into spacetime curvature strong enough to form a horizon
requires an effective stress-energy content for the colliding regions
--- the same open requirement flagged for H1 and H4, since GR curves
in response to stress-energy, not motion of space considered as
already-empty geometry. (iii) The nearest existing, if still
speculative, research precedent for ``collision forms a black hole''
is bubble-collision black-hole formation during an early-universe
phase transition, studied in specific inflationary/phase-transition
models -- but that mechanism is tied to the very early universe and a
specific particle-physics scenario, not to expansion generally or to
an ongoing, occasional process. (iv) Rotation (frame dragging) in
real supermassive black holes is well explained as inherited angular
momentum from the matter and smaller black holes that merged and
accreted to form them -- the standard growth picture H5 sets out to
replace -- so H5 would need to show what a vortex-of-colliding-space
origin predicts that spin-from-accretion-and-mergers does not.

\subsection{Hypothesis 6: Galaxies as Matter Settled into a Vortex's Depression, Replacing Dark Matter}
\label{subsec:h6}
\textbf{Statement (rough form).} A galaxy is matter tracing out the
depression in space caused by a central vortex (a supermassive black
hole, per H5). Matter settles into this depression. Stars orbit the
central vortex at speeds set by the strength and reach of the
depression it causes; faster-than-expected orbital speeds at a
galaxy's edge come from the vortex's own strength and nature, and
dark matter is not required to explain them. Different galaxy
morphologies reflect different forms or strengths of the underlying
spacetime turbulence.

\textbf{Formal statement.}
[PLACEHOLDER: needs an explicit $v(r)$ (orbital speed as a function of
radius) predicted from ``vortex strength,'' to compare against the
measured rotation curves this hypothesis is meant to explain]

\textbf{Motivation.} Correctly identifies the real observation: galaxy
rotation curves are measured to stay flat (roughly constant orbital
speed) out to tens of kiloparsecs, far beyond where the visible
(stellar and gas) mass alone, via ordinary $1/r^2$ gravity, predicts a
declining (Keplerian) curve -- this is the original dark-matter
evidence (Rubin \& Ford, 1970s) and a real target for any competing
explanation.

\textbf{What would falsify it / open problems.}
(i) \emph{Scale mismatch.} A supermassive black hole's own
gravitational reach (its ``sphere of influence'') extends only a few
parsecs to a few tens of parsecs from the galactic center -- it
dominates stellar orbits only there. Flat rotation curves are
measured at tens of \emph{kilo}parsecs, roughly three to four orders
of magnitude farther out, where the central black hole's own
gravitational pull, by standard $1/r^2$ falloff, is already
negligible. For H6 to work, ``the vortex'' has to mean something with
far greater mass and reach than the black hole itself -- at which
point it is doing the explanatory work dark matter currently does,
under a new name, and needs its own mass profile stated and checked
against the data the way a dark-matter halo profile is. (ii) \emph{
Independent evidence beyond rotation curves.} Dark matter is inferred
through several lines of evidence besides rotation curves, each of
which a replacement needs to address: the Bullet Cluster, where
gravitational lensing maps show the dominant mass concentration
spatially separated from the hot, X-ray-emitting gas that holds most
of the ordinary matter in the colliding clusters -- direct evidence
for a gravitating component not tied to where the normal matter sits;
the relative heights of the cosmic microwave background's acoustic
peaks, which fix the ratio of ordinary to total matter in the very
early universe, long before galaxies or their central black holes
existed; and the growth rate of large-scale structure, which needs
the extra gravitational pull of a non-baryonic component to form
galaxies within the age of the universe at all. A vortex tied to a
galaxy's own central black hole has no obvious bearing on any of
these three, since two of them predate galaxy formation and one
(Bullet Cluster) is at cluster, not single-galaxy, scale. (iii) H6's
closest existing competitor is Modified Newtonian Dynamics (MOND;
Milgrom, 1983), which changes the force law itself at low
accelerations and reproduces many individual galaxies' rotation
curves strikingly well without extra matter -- but is also known to
struggle with exactly the cluster-scale and CMB evidence in (ii). If
H6 is a vortex-flavored version of the same idea, it inherits MOND's
real successes but should also expect, and address, MOND's known
difficulties rather than treating rotation curves as the only test.
(iv) \emph{Galaxy-to-galaxy diversity.} Galaxies of similar visible
(and presumably similar central-black-hole) mass are observed to
need very different amounts of inferred dark matter -- some
ultra-diffuse galaxies need a large extra-mass fraction, while a few
reported galaxies appear to need very little. A mechanism tied
tightly to a single central vortex's mass needs to explain this
spread; the standard dark-matter picture explains it as galaxies
having assembled with different amounts of dark-matter halo relative
to their visible mass, through different formation histories.

\subsection{Taxonomy of Spacetime Inconsistencies}
\label{subsec:taxonomy}
H3 through H6 each reach for the same underlying idea -- that
spacetime is not smooth, and that departures from smoothness come in
different forms -- under different names (a tear, a depression, a
vortex, a crack or channel). Treating these as named \emph{species} of
one genus, ``spacetime inconsistency,'' is a workable move, but only
if each species is given its own criteria up front. Otherwise the
genus absorbs any new observation after the fact without having
predicted it beforehand, which trades away falsifiability rather than
gaining explanatory power. The table below is a first attempt at
those criteria; every row currently has open entries that need a
decision before the taxonomy does real work.

\begin{center}
\footnotesize
\begin{tabular}{@{}>{\raggedright\arraybackslash}p{1.6cm}>{\raggedright\arraybackslash}p{2.6cm}>{\raggedright\arraybackslash}p{1.3cm}>{\raggedright\arraybackslash}p{2.3cm}>{\raggedright\arraybackslash}p{1.6cm}>{\raggedright\arraybackslash}p{2.9cm}@{}}
\hline
Species & Used by & Horizon? & Mass-energy source? & Scale & What would test it \\ \hline
Sphere (classical BH) & H3, H4, H5, H6 (galactic core) & Yes & Yes -- standard GR collapse or accretion & Stellar to $10^{10}\,M_\odot$ & Already tested: matches GR/Kerr predictions, event-horizon imaging (EHT) \\
Vortex & H5 (SMBH origin), H6 (galaxy-scale reach) & [PLACEHOLDER: decide] & [PLACEHOLDER: what sources the rotation?] & Galactic ($\sim$kpc--Mpc) & Must reproduce flat rotation curves \emph{and} Bullet Cluster, CMB peaks (\S\ref{subsec:h6}) \\
Depression (micro) & H4 & No (by H4's own revision, \S\ref{subsec:h4} discussion) & [PLACEHOLDER: ordinary local mass, or spontaneous?] & Sub-planetary to planetary & Must not exceed Hawking-evaporation or runaway-accretion bounds for its assumed mass (\S\ref{subsec:h4}) \\
Crack / channel & H6 (cosmic web framing) & [PLACEHOLDER: decide] & [PLACEHOLDER: decide] & $\sim$10--100 Mpc (filaments) & Must reproduce filament/void statistics already matched by gravity + cold dark matter (Zel'dovich approximation; Millennium/IllustrisTNG simulations) \\ \hline
\end{tabular}
\end{center}

The two columns most worth settling first are \emph{Horizon?} and
\emph{Mass-energy source?}, since they are what let H3--H6 borrow (or
not) from established black-hole physics without equivocating between
species: a species with a horizon inherits Hawking radiation and
accretion-growth constraints and must be checked against them, as the
Sphere row already is; a species without one cannot claim to be
``technically a black hole'' when that label is convenient elsewhere
in the paper. The \emph{Mass-energy source} column is the same open
question raised for H1 (\S\ref{subsec:h1}) and H2
(\S\ref{subsec:h2}): under standard GR, curvature is sourced by
stress-energy, so any species claimed to arise without a stated
mass-energy source is proposing an amendment to the field equations
themselves, not a new kind of black hole, and should be argued for at
that level.

% =====================================================================
\section{Consistency Checks}
\label{sec:consistency}
% =====================================================================
% The load-bearing section for scientific credibility. At minimum,
% show (or explicitly flag as future work, if not yet done):

\subsection{Recovery of known limits}
% Does the hypothesis reduce to Newtonian gravity in the weak-field,
% low-velocity limit? Does it reduce to standard GR when [parameter]
% -> [value]? This is usually the single hardest and most important
% requirement for a new gravity-related proposal.
[PLACEHOLDER]

\subsection{Dimensional analysis and units}
% Walk through units for every new quantity introduced; a hypothesis
% that doesn't balance dimensionally is dead on arrival.
[PLACEHOLDER]

\subsection{Known observational constraints}
% Check against, as relevant: equivalence-principle tests (Eotvos-type
% experiments), Solar System tests (Mercury perihelion precession,
% light bending, Shapiro delay), binary pulsar timing, gravitational-wave
% observations (LIGO/Virgo: speed of gravity, waveform shape), and
% cosmological constraints (CMB, large-scale structure) if the
% hypothesis touches cosmological scales. State which constraints
% apply, and whether the hypothesis is consistent with, silent on, or
% in tension with each.
[PLACEHOLDER]

\subsection{Internal consistency}
% Conservation laws, causality (no superluminal signaling unless that
% is itself the explicit, defended claim), symmetry requirements
% (diffeomorphism invariance, Lorentz invariance or a stated and
% motivated departure from it).
[PLACEHOLDER]

% =====================================================================
\section{Predictions}
\label{sec:predictions}
% =====================================================================
% The section that turns a hypothesis into science: at least one
% prediction that differs, even slightly, from standard GR/QFT in a
% way that is in principle measurable, naming the regime (lab-scale,
% Solar-System-scale, astrophysical, cosmological) and roughly what
% precision would be needed to see the difference.
[PLACEHOLDER]

% =====================================================================
\section{Discussion}
\label{sec:discussion}
% =====================================================================
\subsection{What this paper does and does not claim}
[PLACEHOLDER -- restate scope narrowly, as in the Introduction]

\subsection{Open problems and known weaknesses}
% The honesty section. Name the strongest objection to each hypothesis
% and whether it is answered, partially answered, or still open.
[PLACEHOLDER]

\subsection{Relation to Pattern Arithmetic (if applicable)}
% If this paper reuses machinery, vocabulary, or structural moves from
% the Pattern Arithmetic framework, state precisely which ones, and
% whether the connection is load-bearing or merely a shared notational
% habit (e.g. typed primitive objects, explicit verb/operation tables,
% named refusals). Be as careful here as anywhere else about not
% overclaiming a structural analogy as a physical identity.
[PLACEHOLDER or: Not applicable.]

% =====================================================================
\section{Conclusion}
\label{sec:conclusion}
% =====================================================================
[PLACEHOLDER -- 3-5 sentences: what was proposed, how strongly grounded
it currently is, and what the next paper or next section needs to show]

% =====================================================================
\appendix
\section{Derivations}
% Longer derivations that would interrupt the main text.
[PLACEHOLDER or: None yet.]

\section{Glossary of Notation}
% One line per symbol introduced in this paper.
[PLACEHOLDER]

% =====================================================================
% BIBLIOGRAPHY -- replace with \bibliography{} + .bib file, or inline
% \bibitem entries, once sources are chosen. Suggested anchor sources
% for a space/time/gravity paper: Einstein (1915/1916) field equations;
% a standard GR textbook (Misner-Thorne-Wheeler, or Wald, or Carroll);
% the specific observational papers for any constraint cited in
% \S\ref{sec:consistency} (e.g. LIGO/Virgo detection papers, Shapiro
% delay measurements, Eotvos-type equivalence-principle tests).
% =====================================================================
\begin{thebibliography}{9}
\bibitem{PLACEHOLDER1} [PLACEHOLDER citation]
\bibitem{PLACEHOLDER2} [PLACEHOLDER citation]
\end{thebibliography}

\end{document}
